\begin{figure}[p]
\centering
\begin{tikzpicture}[
  x=1cm,y=1cm,
  box/.style={draw=black!65,rounded corners=2pt,fill=black!2,
    align=center,font=\small,inner sep=6pt,text width=6.0cm},
  wide/.style={box,text width=12.7cm},
  arrow/.style={-{Stealth[length=2mm]},line width=.65pt,draw=black!75}
]
\node[wide] (bad) at (0,0)
  {Hypothetical counterexamples (Section~\ref{sec:reduction-and-notation}):\\
   $n\to\infty$, $C_n=N/2^n\to\infty$, $N\le n2^n$};
\node[wide] (small) at (0,-1.8)
  {Weighted densities are close to $1/2$ at small widths\\
   after excluding configurations that yield a cube\\
   (Sections~\ref{sec:bias-versus-purity-at-power-widths}--\ref{sec:a-small-power-key-grid})};
\node[box] (intermediate) at (-3.4,-4.0)
  {Density bounds extend to unequal\\
   widths on the two host sides\\
   (Sections~\ref{sec:asymmetric-pure-patches-under-2}--\ref{sec:intermediate-powers-of-bias})};
\node[box] (clusters) at (3.4,-4.0)
  {Full-dimensional cluster patches\\
   are excluded
   (Section~\ref{sec:full-dimensional-cluster-patches})};
\node[wide] (deep) at (0,-6.3)
  {Discrepancy holds at a small linear width on either side\\
   and a power width on the other\\
   (Sections~\ref{sec:jump-to-large-bias-only-at-linear-budget}--\ref{sec:the-two-deep-orientations-some-estimates})};
\node[wide] (tiling) at (0,-8.3)
  {Disjoint host patches of one color and orientation\\
   receive subcubes specified by initial coordinates\\
   (Sections~\ref{sec:patch-tiling-under-9}--\ref{sec:internal-solver-and-profiles-for-these-tilings})};
\node[box] (high) at (-3.4,-10.8)
  {Larger gains (Section~\ref{sec:high-modes-of-the-tiling})\\
   Odd labels are assigned injectively\\
   Even column sums give Hall};
\node[box] (low) at (3.4,-10.8)
  {Smaller gains (Sections~\ref{sec:low-modes-geometry-and-cell-states}--\ref{sec:late-dynamics-and-its-transfer-problem})\\
   Preserve lists through late assignments\\
   Joint pair bounds give Hall};
\node[wide] (cube) at (0,-13.3)
  {A monochromatic copy of $Q_n$: contradiction\\
   Hence $\sup_{n\ge0}R(Q_n)/2^n<\infty$};
\draw[arrow] (bad) -- (small);
\draw[arrow] (small.south) -- (0,-2.75) -| (intermediate.north);
\draw[arrow] (small.south) -- (0,-2.75) -| (clusters.north);
\draw[arrow] (intermediate.south) |- (0,-5.45) -- (deep.north);
\draw[arrow] (clusters.south) |- (0,-5.45) -- (deep.north);
\draw[arrow] (deep) -- (tiling);
\draw[arrow] (clusters.east) -- ++(.42,0) |- (tiling.east);
\draw[arrow] (tiling.south) -- (0,-9.3) -| (high.north);
\draw[arrow] (tiling.south) -- (0,-9.3) -| (low.north);
\draw[arrow] (high.south) |- (0,-12.05) -- (cube.north);
\draw[arrow] (low.south) |- (0,-12.05) -- (cube.north);
\end{tikzpicture}
\caption{Logical dependencies of the proof. Each reduction is carried out
along the hypothetical counterexample sequence, passing to subsequences
and removing a vanishing fraction of labels where needed. The separate
arrow from cluster exclusion to patch extraction records its second use:
it bounds the measured scales of the remaining cluster patches.
One common color and orientation are fixed before cube allocation.
The final estimates apply to every sequence $C_n\to\infty$, with no
prescribed rate of divergence.}
\label{fig:proof-roadmap}
\end{figure}
